% Section 1: Motivation and significance
\section{Motivation and significance}
\label{sec:introduction}

Forest operations planners make short-horizon decisions about machine allocation,
sequence, and timing under binding logistical constraints. These decisions are evaluated against cost, safety, environmental performance, and reliability of delivery, while operating conditions vary between tenure, terrain, and harvest systems
\cite{henkelman1978gpssv,weintraub1996new,heinimann2007forest};
decision-support systems have long treated these settings as optimization problems
\cite{buongiorno2003decision,kangas2008decision}.
In practice, organizations still combine spreadsheets, geographic information system (GIS) workflows, and local scripts for day-to-day scheduling,  limiting comparability and slowing model transfer
\cite{ulvdal2022uncertainty,nelson2003forestmodels}.

The systematic review by Jaffray et al.~provides the clearest current synthesis of this problem space \cite{jaffray2025systematic}. Among 23 peer-reviewed studies published between 2005 and 2024, the review reports strong technical contributions but recurring structural limitations: high site specificity, limited open artifacts, and weak or no reproducibility support. The design implications of the review are explicit: operational planning tools need standardized inputs, transparent formulations, and reproducible evaluation workflows. These changes in operational planning tools will allow methods, identification of recurring limiting factors and spatial impacts to be compared without rebuilding full software stacks.

This SoftwareX paper is the middle contribution in a review $\rightarrow$ platform $\rightarrow$ modelling sequence. The review defines the gap and requirements; this paper documents the Forest Harvesting Operations Planning System (FHOPS) as reusable infrastructure that operationalizes those requirements and produces a reproducible baseline evidence package; follow-up modelling work then deploys FHOPS on a number of case studies to explore specific design questions (e.g., rolling-horizon configuration) on that shared baseline.

FHOPS combines a scenario contract for operational inputs, heuristic and mixed-integer exact solver pathways, and a shared evaluation layer for runtime, schedule quality, robustness, and costing interpretation. In this context, FHOPS executes machine scheduling in forest operations as a structured, scenario-driven optimization problem with explicit resource, sequencing and timing decisions. The contribution is not a new optimization theorem; it is a reusable, inspectable software that enables equivalent experiments on shared structures with traceable outputs.

In this paper, we make three contributions.
\begin{enumerate}
  \item \textbf{Operational planning software architecture.} We define a domain-first software structure around blocks, machines, shifts, landings, and mobilization rules allowing scheduling analyses that are transferrable across cases with minimal recoding.
  \item \textbf{Canonical formulation + traceability.} We present the operational mixed-integer linear program (MILP) in a canonical equation block with direct links to implementation and regression tests, enabling reviewer-level verification from manuscript equation to executable model.
  \item \textbf{Reproducible validation baseline.} We report benchmark, tuning, playback-robustness, and costing outputs on a public British Columbia (BC) reference ladder (tiny7, small21, med42), with synthetic scaling support, so future modelling studies can reuse a stable baseline.
\end{enumerate}

Scope is intentionally bounded. We focus on software platform capability and reproducible baseline evidence, not on broad claims of external deployment impact. Real-world multi-context validation (including additional system families and rolling-horizon case analyses) is treated as downstream modelling work that reuses the same platform and artifacts rather than redefining them.

Section~\ref{sec:software-description} describes entities, methods, and formulation traceability. Section~\ref{sec:illustrative-example} reports validation evidence around planner-facing operational questions. Section~\ref{sec:impact} discusses deployment relevance, maturity, and transfer limits. Section~\ref{sec:conclusions} closes with claims, limits, and the handoff to companion modelling studies.
